% !PRINTED 685
\weberpage{685}
\S\ 196. \gesperrt{Genus of the field $\Omega$.}

Since in a polygon quotient numerator and denominator are of the same order, it follows from the last formula [\S\ 195, (12)]:
\weq{1}{w_\alpha - 2a = w_\beta - 2b.}

If, then, we denote by $n$ the order and by $w$ the branching number for an arbitrary variable, which, as we have seen in \S\ 193, is an even number, then the integer:
\weq{2}{p = \tfrac{1}{2}w - n + 1,}
is an invariant integer belonging to the field $\Omega$, and it is called the \gesperrt{genus of the field}. That this number cannot be negative follows from \S\ 193, (5) and (9), according to which
\weq{3}{p = (r_2-1)+(r_3-1)+\cdots+(r_{n_1}-1)}
% [sic] German source prints final subscript $n_1$; formulas §193 (5), (9) require $n$.
holds, so that it consists entirely of summands of which none is negative.

Two functions $\alpha$, $\beta$ we call a \gesperrt{primitive pair} of the field $\Omega$ if all functions in $\Omega$ can be expressed rationally by $\alpha$ and $\beta$. If $\alpha$ is of order $m$ and $\beta$ of order $n$, then by \S\ 171 it is necessary and sufficient for a primitive pair that the irreducible equation subsisting between $\alpha$ and $\beta$:
\weq{4}{F(\overset{n}{\alpha},\overset{m}{\beta}) = 0}
be of degree $n$ in $\alpha$ and of degree $m$ in $\beta$.

The functions $F'(\alpha)$, $F'(\beta)$ are now to be represented by polygon quotients. We set:
\weq{5}{\alpha = \dfrac{\mathfrak{M}}{\mathfrak{A}}, \quad \beta = \dfrac{\mathfrak{N}}{\mathfrak{B}},}
in which $\mathfrak{A}$, $\mathfrak{M}$ are of order $m$, and $\mathfrak{B}$, $\mathfrak{N}$ of order $n$. $\mathfrak{A}$ then has no point in common with $\mathfrak{M}$, nor $\mathfrak{B}$ with $\mathfrak{N}$. We set:
\weq{6}{\begin{aligned}
F(\alpha,\beta) &= a_0\alpha^n + a_1\alpha^{n-1} + \cdots + a_n,\\
F'(\alpha) &= na_0\alpha^{n-1} + (n-1)a_1\alpha^{n-2} + \cdots + a_{n-1},\\
\alpha F'(\alpha) &= -a_1\alpha^{n-1} - 2a_2\alpha^{n-2} - \cdots - na_n.
\end{aligned}}

From the second of these equations one concludes that in the denominator of $F'(\alpha)$ there occur no points that are not contained in $\mathfrak{A}$ or in $\mathfrak{B}$, and from the third it follows that this denominator can at most be $= \mathfrak{A}^{n-2}\mathfrak{B}^m$. We therefore set:
\weq{7}{F'(\alpha) = \dfrac{\mathfrak{L}}{\mathfrak{A}^{n-2}\mathfrak{B}^m}.}
